% tikz-tensors-tree.tex -- the tree layout: a stack (tikz-tensors-stack.tex)
% drawn as a tree is usually drawn. It is \tnstack[tree], and everything it
% does not say here it does as a stack.
%
% It differs in three things. Its pitch across is half a stack's, because a
% tree's leaves are tensors of one site each and nothing is set across in its
% slot. A tensor across several sites keeps its own size, centred over them,
% instead of stretching to their width. And an index meets a tensor where the
% tensor's base can take it, so an index between two tensors that do not line
% up turns a corner rather than stretching either -- down from the upper one,
% across half way, and down into the lower one.
\tn@kind{tree}{stack}

\def\tn@tree@geometry#1{\tn@set{#1@pitch}{\tn@treepitch}\tn@set{#1@rise}{\tn@treerise}}
\def\tn@tree@spanwidth{\def\tn@span@w{1}}

% Where an index at x = \tn@cx meets tensor #1, as \tn@xc: straight on if its
% base reaches that far, else as far out as its base comfortably takes an
% index -- three fifths of the way to its edge. A triangle is half as wide at
% mid-height, where `east' is, as at its base.
\def\tn@meet#1{%
  \tn@anc{#1}{center}\let\tn@mx\tn@ax
  \tn@anc{#1}{east}\tn@len\tn@hw{\tn@ax-\tn@mx}%
  \edef\tn@one{\tn@get{n@#1@one}}%
  \ifx\tn@one\@empty\else \tn@len\tn@hw{2*\tn@hw}\fi
  \tn@len\tn@off{abs(\tn@cx-\tn@mx)}%
  \ifdim\tn@off>\tn@hw
    \tn@len\tn@xc{min(max(\tn@cx,\tn@mx-0.6*\tn@hw),\tn@mx+0.6*\tn@hw)}%
  \else \let\tn@xc\tn@cx \fi}
\def\tn@tree@vport#1#2{\tn@stack@vport{#1}{#2}\tn@meet{#1}\let\tn@px\tn@xc}
